\documentclass[../DoAn.tex]{subfiles}
\begin{document}

Phụ lục này mô tả chi tiết cách sinh các tập dữ liệu thực nghiệm và tham số cấu hình của các thuật toán baseline dùng để so sánh, nhằm bảo đảm khả năng tái lập.

\section{Chi tiết sinh tập dữ liệu}
\label{sec:app-data}

Mỗi kịch bản gồm một mạng vật lý đa miền và hai tập VNR (huấn luyện và kiểm thử) tách biệt. Mạng vật lý được chia thành $10$ miền; các nút và liên kết được sinh ngẫu nhiên trong các khoảng tham số ở Bảng~\ref{tab:app-substrate}. Các liên kết liên miền nối các nút biên giữa những miền lân cận, tạo thành một tô-pô liên miền liên thông.

\begin{table}[H]
\centering
\caption{Khoảng tham số sinh mạng vật lý (kịch bản 50 và 100 nút)}
\label{tab:app-substrate}
\begin{tabular}{lcc}
\toprule
\textbf{Tham số} & \textbf{50 nút} & \textbf{100 nút} \\
\midrule
Số miền                & 10 & 10 \\
Số nút                 & 50 & 100 \\
Số liên kết            & 50 & 110 \\
Dung lượng CPU         & 10.5--29.4 & 10.5--29.4 \\
Đơn giá CPU            & 1.1--3.95  & 1.03--4.39 \\
Băng thông liên kết    & 42--93     & 42--96 \\
Đơn giá băng thông     & 0.18--0.99 & 0.29--0.99 \\
Độ trễ xử lý           & 0.25--1.87 & 0.25--1.63 \\
Độ trễ truyền dẫn      & 1.05--9.76 & 2.47--9.76 \\
\bottomrule
\end{tabular}
\end{table}

Các yêu cầu mạng ảo được sinh với tham số ở Bảng~\ref{tab:app-vnr}. Thời điểm đến tuân theo tiến trình Poisson và thời gian sống được sinh ngẫu nhiên, mô phỏng kịch bản trực tuyến nơi VNR đến và rời đi liên tục theo thời gian. Tập huấn luyện gồm $10{,}000$ VNR, tập kiểm thử gồm $3{,}000$ VNR.

\begin{table}[H]
\centering
\caption{Khoảng tham số sinh yêu cầu mạng ảo}
\label{tab:app-vnr}
\begin{tabular}{lc}
\toprule
\textbf{Tham số} & \textbf{Giá trị} \\
\midrule
Số nút mỗi VNR           & 3--7 (trung bình 5) \\
Nhu cầu CPU mỗi nút ảo   & 1.0--8.0 \\
Nhu cầu băng thông mỗi liên kết ảo & 5.0--25.0 \\
Phân phối thời điểm đến  & Poisson \\
Thời gian sống           & ngẫu nhiên \\
Số VNR huấn luyện        & 10{,}000 \\
Số VNR kiểm thử          & 3{,}000 \\
\bottomrule
\end{tabular}
\end{table}

\section{Tham số các thuật toán so sánh}
\label{sec:app-baselines}

Bảng~\ref{tab:app-baseline-params} liệt kê tham số của các thuật toán baseline. \textbf{MP-VNE} dùng cấu hình PSO mạnh hơn ($20$ hạt, $15$ vòng lặp), đóng vai trò trần hiệu năng tham chiếu. \textbf{MP-VNE-V4} dùng đúng tham số theo bài báo gốc ($10$ hạt, $50$ vòng lặp, hệ số quán tính/nhận thức/xã hội $0{,}3/0{,}3/0{,}4$), là baseline được so sánh trực tiếp ở mọi quy mô. Cả hai đều xếp hạng ứng viên theo chi phí ước lượng (PreCost) và ánh xạ liên kết bằng đường đi ngắn nhất.

\begin{table}[H]
\centering
\caption{Tham số PSO của các baseline heuristic}
\label{tab:app-baseline-params}
\begin{tabular}{lcc}
\toprule
\textbf{Tham số} & \textbf{MP-VNE} & \textbf{MP-VNE-V4} \\
\midrule
Số hạt (particles)        & 20  & 10 \\
Số vòng lặp (iterations)  & 15  & 50 \\
Hệ số quán tính $\alpha$  & 0.7 & 0.3 \\
Hệ số nhận thức $\beta$   & 1.5 & 0.3 \\
Hệ số xã hội $\gamma$     & 1.5 & 0.4 \\
Xác suất đột biến         & 0.1 & 0.1 \\
Số ứng viên (top-K)       & 5   & top-K theo miền \\
\bottomrule
\end{tabular}
\end{table}

Mô hình baseline \textbf{R2} (\texttt{il\_mp\_vne\_v17}) là mạng chính sách chỉ qua giai đoạn học bắt chước, chưa tinh chỉnh tăng cường, dùng để cô lập đóng góp của giai đoạn học tăng cường như phân tích ở Mục~\ref{subsec:abl-rl}.

\end{document}
